\clearpage
% Experimental scope and values: authoritative report at 4d6a4a2.
\section{Thiết kế thực nghiệm và kết quả}
\subsection{Giao thức đánh giá chung}
\subsubsection{Đối chứng, cấu hình và chia tập}
Các thực nghiệm đánh giá nhánh RAG theo cấu hình của từng lần chạy, không bao gồm chat dự phòng.
\begin{itemize}
\item \textbf{Đối chứng và chia tập:} Giai đoạn 1 chọn truy xuất/reranker trên tập phát triển (development), không sinh hoặc chấm đáp án. Giai đoạn 2 giữ truy xuất, đối chứng P0-D5; sàng lọc (screening) dùng 80 câu, chấm RAGAS 20 câu, chọn một lần bằng seed 42. Ứng viên cuối (finalist) dùng đủ 281 câu rồi khóa cấu hình trước held-out.
\item \textbf{Mô hình sinh và chấm:} Gemini 3.1 Flash-Lite sinh đáp án với temperature 0, tối đa 512 token, reasoning minimal. GLM-5.3-Flash qua Fireworks chấm năm thước đo RAGAS với reasoning low, tối đa 2.048 token và embedding all-MiniLM-L6-v2.
\item \textbf{Biến kiểm soát:} khi đổi prompt hoặc số đoạn ngữ cảnh (context depth), dùng lại kết quả truy xuất và thứ tự xếp hạng theo mã câu hỏi để cô lập tác động tầng sinh.
\end{itemize}
Cấu hình ứng dụng quy định xử lý yêu cầu; cấu hình thực nghiệm xác định dữ liệu, chỉ mục, tập và biến khảo sát. Hồ sơ lần chạy ghi mô hình, mã băm prompt/trace, danh sách câu hỏi và nhật ký thử lại.
Trong Bảng~\ref{tab:experiment-map}, D là development 281 câu; S là sàng lọc 80 câu với 20 câu chấm RAGAS; H là generation held-out 284 câu; F là retrieval final-test 871 câu. D5 nghĩa là dùng 5 đoạn ngữ cảnh.
\begin{table}[!htbp]\setstretch{1.05}\centering\small
\caption{Tóm tắt thiết lập các thực nghiệm chính.}\label{tab:experiment-map}
\begin{tabularx}{\linewidth}{@{}P{1.45cm}Y Y Y@{}}\toprule
\textbf{Mục tiêu} & \textbf{So sánh và đối chứng} & \textbf{Giữ cố định; tập} & \textbf{Thước đo, quy tắc}\\\midrule
Chọn truy xuất & Retriever/reranker; đối chứng không rerank & Corpus, chia đoạn; D & MRR, nDCG, Hit; đủ coverage \\
Chọn zero-shot & P0--P3, D1/D3/D5; đối chứng P0-D5 & Truy xuất, generator/judge; S rồi D & AC sau guardrail; token, P95 \\
Đo one-shot & P2-1S so với P2-D5 & Truy xuất, D5, generator/judge; D và H & EM/F1, AC, citation; phân tích đánh đổi \\
Đổi generator & Gemini 3.5/3.6 so với 3.1 & P2-1S-D5, truy xuất, judge; D & Chất lượng, token, P95; không chọn lại \\
Kiểm tra & Development so với bài mới & Pipeline đã khóa; F và H & Truy xuất, đáp án; không chọn lại \\\bottomrule
\end{tabularx}\end{table}
\subsubsection{Thước đo và độ bao phủ}
Các thước đo truy xuất và dẫn nguồn dùng tập đoạn bằng chứng chuẩn (gold) ánh xạ từ nhãn dữ liệu; điểm tổng hợp lấy trung bình theo câu.
\begin{table}[H]\centering\small\setstretch{1.05}
\caption{Thước đo và cách diễn giải.}\label{tab:metrics}
\begin{tabularx}{\linewidth}{@{}P{4.5cm}Y@{}}\toprule
\textbf{Thước đo} & \textbf{Nội dung đánh giá}\\\midrule
Hit@$k$ / Recall@$k$ & Hit bằng 1 nếu top $k$ chứa đoạn chuẩn, bằng 0 nếu không; Recall là tỷ lệ đoạn chuẩn tìm được. \\
MRR@$k$ / nDCG@$k$ & MRR là nghịch đảo thứ hạng đoạn chuẩn đầu tiên trong top $k$ (0 nếu không có). nDCG chiết khấu $1/\log_2(r+1)$, chuẩn hóa theo thứ tự lý tưởng. \\
Exact Match / Token F1 & So khớp toàn chuỗi / chồng lấp token; lấy điểm cao nhất trong các đáp án chấp nhận. \\
Citation Validity & Tỷ lệ chỉ số trích dẫn nằm trong phạm vi ngữ cảnh $[1..k]$ (tính hợp lệ cú pháp). \\
Citation Precision / Recall & $|C \cap G| / |C|$ (đoạn dẫn trúng gold) / $|C \cap G| / |G|$ (độ bao phủ gold); $C$ là tập đoạn được dẫn, $G$ là tập đoạn chuẩn. Citation F1 là trung bình điều hòa. \\
Faithfulness / Answer Relevancy & Mức hỗ trợ của ngữ cảnh đối với khẳng định / mức tương tự giữa câu hỏi sinh từ đáp án và câu hỏi gốc~\cite{ragas}. \\
Context Precision / Recall & Mức liên quan và thứ tự các đoạn / mức bao phủ thông tin của đáp án chuẩn~\cite{ragas}. \\
Answer Correctness (AC) & Kết hợp đối chiếu sự kiện và độ tương tự ngữ nghĩa với đáp án chuẩn~\cite{ragasdocs}; điểm chính tầng sinh, không phải tỷ lệ câu đúng. \\\bottomrule
\end{tabularx}\end{table}
\textbf{Tiền xử lý.} Chấm EM/F1 sau khi bỏ chỉ dấu nguồn và nhãn mở đầu, chuyển chữ thường, bỏ dấu câu và chuẩn hóa khoảng trắng. Các điểm dẫn nguồn loại chỉ dấu trùng; câu không có dẫn nguồn nhận điểm 0. Trong bảng, AR là Answer Relevancy; Faith. là Faithfulness; Cit. F1/Valid. là Citation F1/Validity.

\textbf{Phân biệt rào chắn dẫn nguồn.} \textit{Citation Validity} chỉ kiểm tra chỉ số dẫn nguồn có nằm trong dải $[1..k]$ được cung cấp hay không nhằm phát hiện ảo giác cú pháp. Ngược lại, \textit{Citation Precision, Recall và F1} đo lường sự thật khách quan bằng cách đối chiếu mã đoạn được trích dẫn với tập đoạn chuẩn $G$. Đồng thời, \textit{Faithfulness} chỉ đo tính nhất quán nội tại (không bịa thông tin ngoài ngữ cảnh được cấp), trong khi \textit{Citation F1} kiểm tra xem câu trả lời có thực sự neo vào đúng bằng chứng chuẩn hay không.

\textbf{Độ bao phủ (coverage).} Coverage sinh là tỷ lệ câu sinh thành công trên số câu yêu cầu; coverage RAGAS là số câu chấm thành công chia số câu sinh thành công. Tỷ lệ 20/80 ở vòng sàng lọc là phạm vi chấm định trước, không phải lỗi. Cách xử lý bản ghi lỗi, thiếu và tiếp tục lần chạy được nêu ở Phụ lục~\ref{app:reproduction}.
\subsubsection{Bootstrap, guardrail và tài nguyên}
\textbf{Bootstrap.} Sai khác trên cùng câu hỏi là $d_q=m_q(B)-m_q(A)$.
\begin{itemize}
\item \textbf{Giai đoạn 1 (Truy xuất):} 1.000 mẫu bootstrap ghép cặp trên 281 câu development.
\item \textbf{Giai đoạn 2 (Sinh đáp án) và các mở rộng:} 2.000 mẫu theo cụm bài; lấy mẫu 50 bài có hoàn lại, giữ các câu của mỗi bài và tính trung bình theo câu. Cách này giữ sự phụ thuộc giữa các câu cùng nguồn.
\end{itemize}
Khoảng tin cậy 95\% (CI95) lấy các phân vị 2,5\% và 97,5\% của sai khác trung bình. CI chứa 0 chưa đủ để kết luận hai phương án tương đương.

\textbf{Quy tắc chọn ở Giai đoạn 2.} Trên toàn bộ development, áp dụng các ngưỡng bảo vệ chất lượng (guardrail) so với baseline:
\begin{itemize}
\item Coverage sinh và chấm ít nhất 95\%.
\item Faithfulness giảm không quá 0,02; Citation F1/Validity giảm không quá 0,01.
\item Trong các ứng viên đạt, ưu tiên AC; khi chênh lệch nhỏ, lần lượt xét token, độ trễ P95 tầng sinh và độ đơn giản. Guardrail là quy tắc lựa chọn riêng, không thay cho CI.
\end{itemize}

\textbf{Tài nguyên.} P50/P95 là trung vị/phân vị 95, tách truy xuất với gọi mô hình sinh. Token chỉ tính sinh đáp án; điều kiện đo trình bày ở Phụ lục~\ref{app:prompts}.

\FloatBarrier
\subsection{Giai đoạn 1: Lựa chọn cấu hình truy xuất}
\subsubsection{Mục tiêu và thiết lập thực nghiệm}
Mục tiêu là chọn cấu hình đưa bằng chứng lên đầu danh sách trên kho 11.064 bài.
\begin{itemize}
\item \textbf{Sàng lọc retriever:} so sánh bốn dense và bốn sparse; giữ recursive 512/64, lấy 10 ứng viên, chưa rerank. Resolved là biến thể chính; original dùng để phân tích ảnh hưởng của câu hỏi thiếu chủ thể.
\item \textbf{So sánh reranker:} đại diện dense, sparse và hybrid RRF kết hợp ba chế độ: không rerank, MiniLM, BGE-reranker-large. Lấy 20 ứng viên, chấm trên 5 đoạn đầu sau xếp hạng lại.
\item \textbf{Quy tắc chọn:} ưu tiên MRR@5, rồi nDCG@5, Hit@5 và độ trễ; mọi lần chạy phải có đủ bản ghi cho 100\% câu hỏi.
\end{itemize}

% Authoritative results: upstream 4d6a4a2, sections/04_phase1_retrieval.tex.
\subsubsection{Kết quả và phân tích}
\begin{table}[H]\setstretch{1.05}
  \centering
  \caption{Kết quả sàng lọc retriever trên 281 câu resolved.}
  \label{tab:phase1-retriever-screening}
  \small
  \begin{tabular}{lrrrr}
    \toprule
    \textbf{Retriever} & \textbf{MRR@5} & \textbf{nDCG@5} &
    \textbf{Hit@5} & \textbf{P50 (ms)} \\
    \midrule
    BGE-M3 learned sparse          & \textbf{0,8059} & \textbf{0,8317} & \textbf{0,9181} & 76,7 \\
    BM25 Okapi + stemming          & 0,7815 & 0,8123 & 0,9146 & 56,1 \\
    BM25+                          & 0,6881 & 0,7206 & 0,8221 & 101,1 \\
    BM25 Okapi                     & 0,6732 & 0,7087 & 0,8185 & 111,8 \\
    Dense: e5-base-v2              & \textbf{0,6316} & \textbf{0,6684} & \textbf{0,7865} & 15,5 \\
    Dense: bge-small-en-v1.5       & 0,6285 & 0,6589 & 0,7651 & 16,0 \\
    Dense: bge-large-en-v1.5       & 0,6252 & 0,6536 & 0,7544 & 27,4 \\
    Dense: all-MiniLM-L6-v2        & 0,4851 & 0,5165 & 0,6299 & 10,6 \\
    \bottomrule
  \end{tabular}
\end{table}

Trên 281 câu resolved, BGE-M3 đạt điểm cao nhất trong nhóm sparse, còn E5-base-v2 được chọn làm đại diện dense. Các phương pháp sparse vượt trội rõ rệt so với các mô hình dense. Phân tích EDA chỉ ra rằng câu hỏi tin tức neo chặt vào các thực thể định danh rời rạc (tên riêng, địa danh, mốc thời gian với $\text{IDF} \ge 6{,}0$). Khi kho ngữ liệu lớn (11.064 bài), mô hình dense bi-encoder nén thông tin vào vector 768 chiều dễ gặp hiện tượng \textbf{pha loãng thực thể (entity dilution)}, làm mờ ranh giới giữa các sự kiện cùng chủ đề và kéo theo nhiều bài báo gây nhiễu; trong khi sparse dựa vào khớp từ vựng chính xác để cô lập không gian tìm kiếm. Tuy nhiên, CI95 của hiệu nDCG@5 giữa BGE-M3 và BM25 có stemming là $[-0{,}0125;\,0{,}0508]$, nên chưa đủ bằng chứng khẳng định BGE-M3 tốt hơn BM25-stemmed trên quần thể mục tiêu.
\begin{table}[H]\setstretch{1.05}
  \centering
  \caption{MRR@5 và nDCG@5 theo retriever và reranker trên tập development.}
  \label{tab:phase1-reranker-matrix}
  \small
  \setlength{\tabcolsep}{4.5pt}
  \begin{tabular}{lrrr|rrr}
    \toprule
    & \multicolumn{3}{c|}{\textbf{MRR@5}}
    & \multicolumn{3}{c}{\textbf{nDCG@5}} \\
    \textbf{Retriever} & \textbf{Không} & \textbf{MiniLM} & \textbf{BGE-large}
    & \textbf{Không} & \textbf{MiniLM} & \textbf{BGE-large} \\
    \midrule
    BGE-M3 sparse & 0,8059 & 0,8387 & \textbf{0,8797} & 0,8317 & 0,8642 & \textbf{0,8976} \\
    e5-base dense & 0,6316 & 0,7783 & 0,7997 & 0,6684 & 0,7994 & 0,8172 \\
    Hybrid RRF    & 0,7192 & 0,7938 & 0,8203 & 0,7474 & 0,8164 & 0,8405 \\
    \bottomrule
  \end{tabular}
\end{table}

\begin{table}[H]\setstretch{1.05}
  \centering
  \caption{Độ trễ P50 của ma trận truy xuất Giai đoạn 1 trên cùng môi trường chạy.}
  \label{tab:phase1-reranker-latency}
  \small\setlength{\tabcolsep}{6pt}
  \begin{tabular}{lrrr}
    \toprule
    \textbf{Retriever} & \textbf{Không rerank} & \textbf{MiniLM} &
    \textbf{BGE-large} \\
    \midrule
    BGE-M3 sparse & 66,1 ms & 163,9 ms & 510,1 ms \\
    e5-base dense & 17,7 ms & 127,3 ms & 540,5 ms \\
    Hybrid RRF    & 90,4 ms & 193,3 ms & 609,8 ms \\
    \bottomrule
  \end{tabular}
\end{table}

\begin{itemize}
\item \textbf{Reranker:} với BGE-M3, MiniLM nâng MRR@5 từ 0,8059 lên 0,8387; BGE-reranker-large đạt 0,8797. Đổi lại, P50 tăng từ 163,9 ms với MiniLM lên 510,1 ms với BGE-large.
\item \textbf{Hybrid RRF:} không vượt BGE-M3 ở cả ba chế độ xếp hạng lại; phép hợp nhất chưa đem lại lợi ích trong thiết lập khảo sát.
\end{itemize}

\noindent\textbf{Kết luận.} BGE-M3 top 20 kết hợp BGE-reranker-large top 5 được chọn cho tầng sinh, với Hit@5/Recall@5 trên development là 0,9573/0,9555.


\subsection{Giai đoạn 2: Lựa chọn cấu hình sinh câu trả lời}
\subsubsection{Lựa chọn prompt zero-shot và số đoạn ngữ cảnh}
Mục tiêu là cải thiện độ đúng câu trả lời, giữ mức hỗ trợ bởi bằng chứng và chất lượng dẫn nguồn. Truy xuất cố định ở recursive 512/64, BGE-M3 top 20, BGE-reranker-large top 5; generator và judge theo giao thức chung. Bốn prompt khác nhau ở chỉ thị:
\begin{itemize}
\item \textbf{P0 — đối chứng:} trả lời dựa trên ngữ cảnh và dẫn nguồn.
\item \textbf{P1 — bám sát bằng chứng:} nhấn mạnh yêu cầu không suy đoán.
\item \textbf{P2 — dạng đáp án:} trả lời ngắn, đúng loại thông tin được hỏi.
\item \textbf{P3 — chọn bằng chứng:} đối chiếu chủ thể, sự kiện, nơi và thời gian để tránh trộn thông tin.
\end{itemize}

Sàng lọc bốn prompt với 5 đoạn; sau đó thử P0/P2 với 1, 3 hoặc 5 đoạn đầu cùng danh sách rerank, dùng lại kết quả 5 đoạn đã có. P2-D3/D5 được chấm đủ development cùng P0-D5; áp dụng guardrail trước khi chọn theo AC.
% Authoritative results: upstream 4d6a4a2, sections/05_phase2_generation.tex.
% Low-score sample and semantic caveat: docs/reports/phase2/report.md, section 6.
\subsubsection{Đối chứng và sàng lọc}
Baseline dùng P0 với 5 đoạn ngữ cảnh (P0-D5) trên toàn bộ development.
\begin{table}[H]\setstretch{1.05}\centering\small
\caption{Baseline P0-D5 trên 281 câu development.}\label{tab:phase2-baseline}
\begin{tabularx}{\linewidth}{@{}YrYr@{}}\toprule
\textbf{Thước đo} & \textbf{Giá trị} & \textbf{Thước đo} & \textbf{Giá trị}\\\midrule
Exact Match & 0,0000 & Faithfulness & 0,9761 \\
Token F1 & 0,2648 & Context Precision & 0,8980 \\
Answer Correctness & 0,6308 & Context Recall & 0,9644 \\
Answer Relevancy & 0,7950 & Coverage & 100\% \\
Citation F1 & 0,8025 & P50 toàn chuỗi (s) & 2,77 \\
Citation Validity & 0,9893 & P95 toàn chuỗi (s) & 53,63 \\\bottomrule
\end{tabularx}\end{table}
\textbf{Phân tích lỗi.} Trong mẫu 30 câu có AC dưới 0,5, có 24 câu vẫn tìm được đoạn chuẩn trong top 5. Các vấn đề gồm trả lời dài, sai loại thông tin hoặc chọn nhầm chi tiết. Điểm thấp không luôn đồng nghĩa sai ngữ nghĩa; mẫu này không đại diện cho tỷ lệ lỗi toàn tập.
\begin{table}[H]\setstretch{1.05}
  \centering
  \caption{Sàng lọc prompt trên 80 câu; RAGAS dùng cùng 20 câu calibration.}
  \label{tab:phase2-prompt-screening}
  \small
  \begin{tabular}{lrrrrrr}
    \toprule
    \textbf{Prompt} & \textbf{Token F1} & \textbf{AC} & \textbf{Faith.} &
    \textbf{AR} & \textbf{Cit. F1} & \textbf{Cit. Valid.} \\
    \midrule
    P0 & 0,2417 & 0,7118 & 1,0000 & 0,8131 & 0,7808 & 0,9750 \\
    P1 & 0,2570 & 0,6885 & 0,9750 & 0,7426 & 0,7863 & 0,9625 \\
    P2 & \textbf{0,3999} & \textbf{0,8363} & \textbf{1,0000} & 0,7388 & 0,8133 & \textbf{0,9750} \\
    P3 & 0,2790 & 0,7190 & 0,9750 & \textbf{0,8153} & \textbf{0,8146} & 0,9625 \\
    \bottomrule
  \end{tabular}
\end{table}

P1 bị loại vì AC và Faithfulness thấp hơn P0. P3 có AC cao hơn P0, nhưng P2 dẫn đầu về AC và Token F1 nên được khảo sát tiếp cùng P0. RAGAS chỉ chấm cùng 20 câu calibration; đáp án và dẫn nguồn chấm trên 80 câu.
\begin{table}[H]\setstretch{1.05}
  \centering
  \caption{Ảnh hưởng của số đoạn ngữ cảnh trên tập sàng lọc.}
  \label{tab:phase2-depth-screening}
  \small
  \begin{tabular}{llrrrrrr}
    \toprule
    \textbf{Prompt} & \textbf{Depth} & \textbf{Token F1} & \textbf{AC} &
    \textbf{Faith.} & \textbf{AR} & \textbf{Cit. F1} & \textbf{Cit. Valid.} \\
    \midrule
    P0 & 1 & 0,2543 & 0,7063 & 0,9750 & 0,7365 & 0,7250 & 0,8167 \\
    P0 & 3 & 0,2527 & 0,7046 & 0,9875 & 0,7667 & 0,7792 & 0,9500 \\
    P0 & 5 & 0,2417 & 0,7118 & 1,0000 & 0,8131 & 0,7808 & 0,9750 \\
    \midrule
    P2 & 1 & \textbf{0,5809} & 0,8671 & 0,9500 & 0,4825 & 0,5500 & 0,6125 \\
    P2 & 3 & 0,5311 & \textbf{0,8795} & 0,9250 & 0,6806 & 0,8125 & 0,9500 \\
    P2 & 5 & 0,3999 & 0,8363 & \textbf{1,0000} & 0,7388 & \textbf{0,8133} & \textbf{0,9750} \\
    \bottomrule
  \end{tabular}
\end{table}

Trên tập sàng lọc, P2-D1 và P2-D3 có AC cao hơn P2-D5, nhưng D1 giảm mạnh Citation F1/Validity. Nhóm đưa P2-D3 và P2-D5 vào vòng cuối trên development.
\subsubsection{Kết quả và lựa chọn}
\begin{table}[H]\setstretch{1.05}
  \centering
  \caption{Kết quả finalist trên toàn bộ development.}
  \label{tab:phase2-zero-shot-finalists}
  \small
  \begin{tabular}{lrrrrrr}
    \toprule
    \textbf{Cấu hình} & \textbf{Token F1} & \textbf{AC} & \textbf{Faith.} &
    \textbf{AR} & \textbf{Cit. F1} & \textbf{Cit. Valid.} \\
    \midrule
    P0-D5 & 0,2648 & 0,6308 & 0,9761 & \textbf{0,7950} & 0,8025 & \textbf{0,9893} \\
    P2-D3 & \textbf{0,5549} & \textbf{0,7711} & 0,9561 & 0,5694 & 0,8381 & 0,9751 \\
    P2-D5 & 0,4191 & 0,7350 & \textbf{0,9801} & 0,7092 & \textbf{0,8391} & 0,9858 \\
    \bottomrule
  \end{tabular}
\end{table}

\begin{itemize}
\item \textbf{P2-D3:} AC cao nhất nhưng không qua guardrail. Citation Validity giảm 0,0142 so với P0-D5, vượt ngưỡng 0,01; Faithfulness giảm 0,020029, vượt nhẹ ngưỡng 0,02.
\item \textbf{P2-D5:} AC tăng 0,1043, CI95 $[0{,}0844;\,0{,}1240]$; Citation F1 tăng 0,0366 và Faithfulness trung bình cao hơn đối chứng.
\end{itemize}

\noindent\textbf{Kết luận.} P2-D5 đạt đủ guardrail và được khóa trước held-out. Chỉ thị trả đúng loại thông tin cải thiện chất lượng trong thiết lập này; AC tăng riêng lẻ chưa đủ nếu dẫn nguồn giảm quá ngưỡng.


\subsubsection{Thử nghiệm one-shot prompt}
Theo góp ý của giảng viên hướng dẫn, nhóm tiến hành thực nghiệm kiểm tra liệu việc bổ sung một ví dụ mẫu (one-shot prompt) có giúp mô hình định hình đáp án ngắn và đặt vị trí dẫn nguồn chính xác hơn hay không. Prompt P2-1S bổ sung một ví dụ tổng hợp về Northbridge/Riverside vào P2, hoàn toàn không lấy từ corpus hoặc tập đánh giá (Phụ lục~\ref{app:prompts}). Hai phương án dùng cùng truy xuất, depth 5, Gemini 3.1 Flash-Lite và judge; đối chứng P2-D5 được tái sử dụng, không sinh hoặc chấm lại.
% Original finalist P2-D5 scores are the frozen baseline; one-shot scores are paired by question.
% Raw deterministic_scores.jsonl confirms EM=26/281; the source report printed 0 in error.
\begin{table}[H]\setstretch{1.05}\centering\small
\caption[Zero-shot và one-shot: điểm số và sai khác ghép cặp.]{Zero-shot và one-shot: development 281 câu, held-out 284 câu. $\Delta$ và CI95 so sánh P2-1S với P2-D5 finalist trên development.}\label{tab:result-phase2d}
\begin{tabular}{@{}lrrrlr@{}}\toprule
\textbf{Thước đo} & \textbf{\shortstack{P2\\dev}} & \textbf{\shortstack{P2-1S\\dev}} & \textbf{$\Delta$} & \textbf{CI95} & \textbf{\shortstack{P2-1S\\held-out}}\\\midrule
Exact Match & 0,0925 & 0,3559 & --- & --- & 0,3239 \\
Token F1 & 0,4191 & 0,6267 & $+0{,}2076$ & $[0{,}1728;\,0{,}2427]$ & 0,6032 \\
Answer Correctness & 0,7350 & 0,8015 & $+0{,}0664$ & $[0{,}0444;\,0{,}0883]$ & 0,7557 \\
Faithfulness & 0,9801 & 0,9603 & $-0{,}0199$ & $[-0{,}0388;\,-0{,}0028]$ & 0,9082 \\
Answer Relevancy & 0,7092 & 0,5251 & $-0{,}1841$ & $[-0{,}2233;\,-0{,}1458]$ & 0,4688 \\
Citation F1 & 0,8391 & 0,8441 & $+0{,}0050$ & $[-0{,}0104;\,0{,}0229]$ & 0,7315 \\
Citation Validity & 0,9858 & 0,9786 & $-0{,}0071$ & $[-0{,}0182;\,0{,}0000]$ & 0,9437 \\
Coverage & 100\% & 100\% & --- & --- & 100\% \\\bottomrule
\end{tabular}\end{table}
\begin{itemize}
\item \textbf{Cải thiện:} AC và Token F1 tăng; CI95 của hai sai khác không chứa 0. Citation F1 chưa có bằng chứng thay đổi.
\item \textbf{Đánh đổi:} Faithfulness và Answer Relevancy giảm. Mức giảm Faithfulness trung bình sát ngưỡng 0,02; one-shot không cải thiện đồng đều các tiêu chí.
\end{itemize}

% Length: word regex after removing all numeric citation groups (including [1, 2]) from raw answers.
\textbf{Phân tích.}
\begin{itemize}
\item \textbf{Quan sát:} trung vị độ dài đáp án (bỏ chỉ dấu nguồn) giảm từ 11 xuống 7 từ; số đáp án không quá ba từ tăng từ 23 lên 105. Số câu có Relevancy dưới 0,5 tăng từ 54 lên 133; 109/133 câu one-shot này vẫn đạt AC ít nhất 0,8.
\item \textbf{Ví dụ:} với câu hỏi về nơi tổ chức giải Swamp Soccer năm 2009, đáp án đầy đủ đạt Relevancy 1,00; ``Scotland [1]'' đạt 0,037 dù Correctness bằng 1,00.
\end{itemize}
RAGAS sinh câu hỏi ngược từ đáp án để tính Relevancy. Đáp án quá ngắn có thể mất chủ thể và sự kiện; cách giải thích này phù hợp với ví dụ, chưa được kiểm chứng trên toàn bộ các trường hợp giảm điểm.

Trên held-out, các thước đo chất lượng của one-shot đều thấp hơn development.

\noindent\textbf{Kết luận.} Trên development, mức giảm Faithfulness và các điểm dẫn nguồn vẫn trong guardrail theo điểm trung bình. Khảo sát diễn ra sau đánh giá held-out zero-shot, ngoài vòng lựa chọn đã khóa. Kết quả phục vụ phân tích đánh đổi, không thay cấu hình P2-D5; việc dùng lại held-out không tạo tập kiểm tra độc lập mới.


\subsubsection{Thử nghiệm so sánh các mô hình sinh}
Theo góp ý của giảng viên hướng dẫn về việc khảo sát tác động khi thay đổi mô hình sinh (generator), thực nghiệm này đánh giá xem việc sử dụng các mô hình khác nhau, đặc biệt là các mô hình mạnh hơn (như Gemini 3.5 Flash-Lite và Gemini 3.6 Flash), mang lại hiệu quả ra sao trên toàn chuỗi RAG, và liệu mức cải thiện về chất lượng câu trả lời có thực sự tương xứng để bù đắp chi phí tài nguyên gia tăng (độ trễ suy luận P95 và lượng token tiêu thụ) hay không. Thử nghiệm giữ nguyên 281 câu development, trace truy xuất, cấu hình P2-1S-D5 và judge; chỉ thay Gemini 3.1 Flash-Lite bằng 3.5 Flash-Lite hoặc 3.6 Flash. Cả ba dùng temperature 0, tối đa 512 token, reasoning minimal. Phép so sánh ghép cặp theo câu dùng kết quả 3.1 đã lưu, không gọi lại đối chứng.
% Authoritative results: upstream 4d6a4a2, sections/05_phase2_generation.tex.
\begin{table}[H]\setstretch{1.05}
  \centering
  \caption{So sánh mô hình sinh với P2-1S-D5 trên development.}
  \label{tab:phase2-generator-ablation}
  \small
  \begin{tabular}{lrrrrrr}
    \toprule
    \textbf{Generator} & \textbf{Token F1} & \textbf{AC} & \textbf{Faith.} &
    \textbf{AR} & \textbf{Cit. F1} & \textbf{Cit. Valid.} \\
    \midrule
    Gemini 3.1 Flash-Lite & 0,6267 & 0,8015 & \textbf{0,9603} & \textbf{0,5251} & \textbf{0,8441} & \textbf{0,9786} \\
    Gemini 3.5 Flash-Lite & 0,6803 & \textbf{0,8164} & 0,9557 & 0,4351 & 0,8351 & 0,9680 \\
    Gemini 3.6 Flash & \textbf{0,6934} & 0,8011 & 0,9472 & 0,3837 & 0,8248 & 0,9644 \\
    \bottomrule
  \end{tabular}
\end{table}

\begin{table}[H]\setstretch{1.05}\centering\small
\caption{Sai khác Answer Correctness và tài nguyên tầng sinh trên 281 câu}\label{tab:generator-operation}
\begin{tabularx}{\linewidth}{@{}YrrP{3.5cm}r@{}}\toprule
\textbf{Generator} & \textbf{$\Delta$ AC} & \textbf{LLM P95 (s)} & \textbf{CI95 của $\Delta$ AC} & \textbf{Token ra}\\\midrule
3.1 Flash-Lite & --- & 1,02 & Đối chứng & 3.847 \\
3.5 Flash-Lite & $+0{,}0149$ & 0,89 & $[-0{,}0052;\,0{,}0359]$ & 3.112 \\
3.6 Flash & $-0{,}0003$ & 3,50 & $[-0{,}0204;\,0{,}0195]$ & 3.249 \\\bottomrule
\end{tabularx}\end{table}
\begin{itemize}
\item \textbf{Gemini 3.5 Flash-Lite:} AC trung bình cao nhất, nhưng CI95 của sai khác so với 3.1 chứa 0.
\item \textbf{Gemini 3.6 Flash:} Token F1 cao nhất nhưng chưa cho thấy cải thiện AC; Citation F1/Validity giảm, LLM P95 cao hơn.
\item \textbf{Gemini 3.1 Flash-Lite:} Faithfulness, Relevancy và hai thước đo dẫn nguồn cao nhất trong ba mô hình khảo sát.
\end{itemize}

Mỗi lần chạy dùng 666.604 token đầu vào; token đầu ra và LLM P95 ở Bảng~\ref{tab:generator-operation}. Token không gồm chấm điểm. Độ trễ phản ánh điều kiện từng lần chạy, không tách riêng ảnh hưởng của mô hình và dịch vụ cung cấp.

\noindent\textbf{Kết luận.} Đổi mô hình tạo đánh đổi về đáp án, bằng chứng, dẫn nguồn và thời gian. Khảo sát chỉ dùng development, không có held-out riêng từng mô hình và không thay cấu hình đã khóa.


\subsection{Đánh giá trên tập kiểm tra độc lập (Held-out Evaluation)}\label{sec:heldout}
Để kiểm định khả năng tổng quát hóa của toàn bộ pipeline RAG đã khóa (BGE-M3 sparse + BGE-reranker-large + Prompt P2-D5), hệ thống được đánh giá độc lập đúng một lần trên tập dữ liệu chưa từng tiếp xúc. Retrieval final-test dùng 871 câu/150 bài chưa dùng để chọn cấu hình; generation held-out dùng 284 câu/50 bài thuộc tập đó. Cấu hình đã khóa được giữ nguyên. Development và held-out tách theo bài, nên không thể so sánh ghép cặp theo câu giữa hai tập. Phân tầng theo đoạn chuẩn trong ngữ cảnh giúp phân biệt hạn chế truy xuất và sử dụng bằng chứng; thiếu đoạn chuẩn không đồng nghĩa hoàn toàn thiếu bằng chứng hợp lý.
% Authoritative final-test results: upstream 4d6a4a2, Phase 1 and Phase 2 sections.
\par\noindent\begin{minipage}[t]{0.49\linewidth}
\begin{table}[H]\setstretch{1.05}
  \centering
  \caption{Truy xuất: development và final-test.}
  \label{tab:phase1-final-test}
  \small
  \begin{tabularx}{\linewidth}{@{}Yrr@{}}
    \toprule
    \textbf{Thước đo} & \textbf{Dev. (281)} &
    \textbf{Final (871)} \\
    \midrule
    MRR@5       & 0,8797 & 0,8112 \\
    nDCG@5      & 0,8976 & 0,8313 \\
    Hit@5       & 0,9573 & 0,8978 \\
    Recall@5    & 0,9555 & 0,8955 \\
    P50 latency & 512,7 ms & 550,0 ms \\
    P95 latency & 531,0 ms & 567,0 ms \\
    Coverage    & 100\% & 100\% \\
    \bottomrule
  \end{tabularx}
\end{table}
\end{minipage}\hfill
\begin{minipage}[t]{0.49\linewidth}
\begin{table}[H]\setstretch{1.05}
  \centering
  \caption{P2-D5 trên development và held-out.}
  \label{tab:phase2-zero-shot-heldout}
  \small
  \begin{tabularx}{\linewidth}{@{}Yrr@{}}
    \toprule
    \textbf{Thước đo} & \textbf{Dev. (281)} & \textbf{Held-out (284)} \\
    \midrule
    Token F1 & 0,4191 & 0,4313 \\
    Answer Correctness & 0,7350 & 0,7157 \\
    Faithfulness & 0,9801 & 0,9249 \\
    Answer Relevancy & 0,7092 & 0,6485 \\
    Citation F1 & 0,8391 & 0,7289 \\
    Citation Validity & 0,9858 & 0,9437 \\
    Coverage & 100\% & 100\% \\
    \bottomrule
  \end{tabularx}
\end{table}
\end{minipage}
\par


\noindent\textbf{Truy xuất.} Pipeline xử lý đủ 871/871 câu: 782 câu có đoạn chuẩn trong top 5, 89 câu có Hit@5 bằng 0. Reranker tăng MRR@5 từ 0,7269 lên 0,8112 và nDCG@5 thêm 0,0758, CI95 $[0{,}0504;\,0{,}1017]$ (2.000 mẫu bootstrap ghép cặp theo 150 bài).


\subsection{Khảo sát phụ và thử nghiệm thăm dò (Ablation Studies)}\label{sec:side-ablations}
Tương tự như báo cáo bảo vệ ban đầu, bên cạnh hai giai đoạn thực nghiệm trọng tâm, nhóm tiến hành hai khảo sát phụ mang tính thăm dò nhằm kiểm chứng tính tối ưu và lý giải vì sao nhóm kiên định bảo toàn cấu hình đã chốt, không mở thêm vòng lựa chọn trên tập held-out:

\subsubsection{Khảo sát phụ 1: Chiến lược phân đoạn văn bản}\label{sec:chunking}
Khảo sát kiểm tra ảnh hưởng của ranh giới đoạn tới bằng chứng và câu trả lời:
\begin{itemize}
\item \textbf{C0 — đối chứng:} recursive 512/64.
\item \textbf{C1 và C2:} C1 ghép các câu, C2 ghép các đoạn văn; cùng giới hạn 512 token.
\item \textbf{C3 — phân cấp:} truy xuất và xếp hạng đoạn con 256/32, mở rộng sang đoạn cha 512/64 rồi loại cha trùng. Citation trỏ tới đoạn cha thực sự gửi cho mô hình sinh.
\end{itemize}

Khảo sát giữ corpus, 281 câu development và cấu hình truy xuất/sinh/chấm của Giai đoạn 2. Vòng sàng lọc đánh giá truy xuất trên 281 câu, sinh trên 80 câu và RAGAS trên 20 câu; sau đó xét độ sâu và finalist trên đủ development. CI dùng 2.000 mẫu bootstrap ghép cặp theo bài, seed 42. Ngoài guardrail của tầng sinh, Hit@5 và Recall@5 không được giảm quá 0,01 so với C0-D5. Tổng ngân sách token ngữ cảnh không cố định giữa các phương án, nên độ dài cũng thuộc tác động được đo.
% Sources: phase2c/notes_thang.md; screening CSV (parent_* for C3);
% paired_significance.json; finalist_lock_c3_d3.json. All under docs/reports/.
\paragraph{Kết quả và phân tích.}
Bảng~\ref{tab:chunking-screen} trình bày kết quả truy xuất của vòng sàng lọc. Với C3, bảng dùng các đoạn cha được chuyển cho mô hình sinh; các chỉ số trên đoạn con không được dùng cho phép so sánh này. C1 và C2 tăng nhẹ các chỉ số, còn C3 thấp hơn C0. C3 vẫn được khảo sát tiếp theo biên sàng lọc 0,02; ngưỡng loại sớm này khác guardrail 0,01 ở vòng cuối.
\begin{table}[H]\centering\small\setstretch{1.05}
\caption{Sàng lọc chia đoạn trên 281 câu, chấm theo ngữ cảnh chuyển cho tầng sinh.}\label{tab:chunking-screen}
\begin{tabular}{@{}lrrrrr@{}}\toprule
\textbf{Phương án} & \textbf{Số đoạn lập chỉ mục} & \textbf{Hit@5} & \textbf{Recall@5} & \textbf{MRR@5} & \textbf{nDCG@5}\\\midrule
C0: recursive & 22.766 & 0,9573 & 0,9537 & 0,8791 & 0,8963\\
C1: theo câu & 22.014 & 0,9680 & 0,9591 & 0,8796 & 0,8979\\
C2: theo đoạn & 22.018 & 0,9609 & 0,9555 & 0,8817 & 0,8989\\
C3: phân cấp & 49.218 & 0,9466 & 0,9431 & 0,8466 & 0,8700\\\bottomrule
\end{tabular}
\end{table}

\textbf{Sàng lọc.}
\begin{itemize}
\item \textbf{Chọn C3 để khảo sát tiếp:} C1--C3 dùng P2-D5 trên 80 câu, cùng 20 câu chấm RAGAS. C3 đạt AC 0,8094 so với C0 là 0,8363; Faithfulness 0,9750 cao hơn C1 (0,9433) và C2 (0,9100). Nhóm tiếp tục C3 vì giữ Faithfulness tốt hơn, dù Citation F1 còn thấp; mẫu này chưa đủ để kết luận C3 tốt hơn.
\item \textbf{Chọn độ sâu:} với C3, D1 bị loại vì Citation F1/Validity chỉ đạt 0,5250/0,6250. D3 có AC cao hơn D5; D5 có Faithfulness và Citation F1 cao hơn D3. Cả D3/D5 vào vòng cuối trên 281 câu.
\end{itemize}

\textbf{Đối chứng.} C0 ở Bảng~\ref{tab:chunking-screen} dùng trace truy xuất của vòng sàng lọc chia đoạn, có Recall@5 bằng 0,9537. Vòng cuối tái sử dụng C0-D3/D5 của Giai đoạn 2, với Recall@5 bằng 0,9555 (Bảng~\ref{tab:phase1-final-test}). Sai khác ở Bảng~\ref{tab:chunking-final} được tính ghép cặp từ điểm từng câu trước khi làm tròn. Hit/Recall@5 chấm danh sách năm đoạn sau mở rộng cha; D3/D5 chỉ giới hạn phần đưa cho mô hình sinh.

\begin{table}[H]\centering\small\setstretch{1.05}
\caption{Ứng viên cuối chia đoạn: sai khác so với C0-P2-D5 trên 281 câu.}\label{tab:chunking-final}
\begin{tabular}{@{}lrrrrr@{}}\toprule
\textbf{Cấu hình} & \textbf{AC} & \textbf{$\Delta$ AC} & \textbf{$\Delta$ Cit. F1} & \textbf{$\Delta$ Hit@5} & \textbf{$\Delta$ Recall@5}\\\midrule
C3-P2-D3 & 0,7730 & $+0{,}0380$ & $-0{,}0284$ & $-0{,}0107$ & $-0{,}0125$\\
C3-P2-D5 & 0,7415 & $+0{,}0065$ & $-0{,}0275$ & $-0{,}0107$ & $-0{,}0125$\\\bottomrule
\end{tabular}
\end{table}
\begin{itemize}
\item \textbf{Quyết định:} cả hai ứng viên đạt coverage 100\% nhưng vi phạm ngưỡng giảm 0,01 của Citation F1, Hit@5 và Recall@5. C0-P2-D5 được giữ nguyên.
\item \textbf{Ảnh hưởng của độ sâu:} mức tăng AC 0,0380 của C3-D3 bao gồm cả thay đổi độ sâu; so với C0-D3 (AC 0,7711), chênh lệch chỉ 0,0019.
\item \textbf{Cơ chế suy giảm của C3:} C3-D3 đạt AC bề mặt cao ($0{,}7730$) nhưng Citation F1 sụt giảm mạnh ($-0{,}0284$), trượt rào chắn an toàn. Nguyên nhân là dù đoạn con 256 token giúp định vị bằng chứng tốt, việc mở rộng trả về đoạn cha 512 token để sinh khiến chỉ dấu trích dẫn \code{[1]} trỏ vào cả đoạn cha rộng lớn, dẫn đến hiện tượng \textbf{pha loãng độ đặc hiệu của nguồn dẫn (citation specificity dilution)} và làm lệch sự đối chiếu với tập đoạn chuẩn gốc.
\item \textbf{Giới hạn:} Citation Validity gần 0,98 chỉ xác nhận chỉ dấu hợp lệ, không bảo đảm đoạn được dẫn khớp bằng chứng chuẩn. Mở rộng cha chưa bảo đảm khớp bằng chứng tốt hơn; phép so sánh chưa tách riêng tác động từng bước.
\end{itemize}


\subsubsection{Khảo sát phụ 2: Chính sách từ chối khi thiếu bằng chứng}\label{sec:abstention}
Khảo sát đánh giá khả năng chủ động từ chối trả lời khi các đoạn được cung cấp không đủ bằng chứng. Bộ 200 tình huống (case) gồm 140 development và 60 final, có các nhóm: đối chứng trả lời được, bỏ bài nguồn, bỏ bằng chứng có kiểm soát, truy xuất thiếu bằng chứng tự nhiên, bằng chứng yếu, câu hỏi phản thực và câu ngoài phạm vi. Hai tập có lần lượt 61 và 26 case trả lời được. Đáp án gốc của case thiếu bằng chứng chỉ mô tả sự kiện nguồn; hành vi mong đợi trên ngữ cảnh được cung cấp là từ chối.

Các biến thể cùng nguồn được giữ trong cùng tập đánh giá. Thử nghiệm dùng kho riêng gồm 19.259 đoạn, không chứa bốn bài nguồn của câu hỏi ngoài phạm vi; không dùng chỉ mục 22.766 đoạn của Giai đoạn 2. Cổng điểm chỉ áp dụng cho 153 case truy xuất toàn kho; 47 case dùng ngữ cảnh đã kiểm soát không qua cổng.

\textbf{Chính sách so sánh.} Generator là Gemini 3.1 Flash-Lite, temperature 0, tối đa 512 token, reasoning minimal; không dùng judge RAGAS.
\begin{itemize}
\item \textbf{B0:} giữ prompt P2 (từ chối tự nhiên).
\item \textbf{B1:} yêu cầu JSON với trường \code{answerability} và đổi chỉ dẫn về cách trả lời.
\item \textbf{B2:} áp cổng điểm xếp hạng lại lên đầu ra B1 đã lưu, không sinh thêm.
\end{itemize}
\textbf{Quy tắc chọn.} Ngưỡng B2 được hiệu chuẩn trên development với false-abstention không quá 10\%. Chính sách được chọn phải thỏa các điều kiện: false-abstention không quá 10\%; tỷ lệ sinh thành công ít nhất 98\%; trên nhóm trả lời được, Token F1 giảm không quá 0,02 và Citation Validity giảm không quá 0,01; ưu tiên false-answer thấp nhất.

\textbf{Hiệu chỉnh bộ đánh giá và CI.} Ba case development được làm rõ sau rà soát; ngưỡng và chính sách được chọn lại, đầu ra final giữ nguyên (Phụ lục~\ref{app:reproduction}). CI dùng 2.000 mẫu bootstrap ghép cặp theo case, seed 42.
% Sources: phase3_final_results.json (post-review amendment),
% policy_comparison.csv, policy_significance.json, docs/reports/phase3/README.md.
\paragraph{Kết quả và phân tích.}
Các thước đo có phạm vi khác nhau:
\begin{itemize}
\item \textbf{Token F1/Citation Validity:} chỉ tính trên nhóm trả lời được.
\item \textbf{False-answer:} lỗi quyết định trên nhóm thiếu bằng chứng (79 case development, 34 case final), gồm trả lời thay vì từ chối và lỗi sinh; không phải tỷ lệ câu trả lời sai sự thật.
\item \textbf{False-abstention:} lỗi quyết định trên nhóm trả lời được, cũng tính lỗi sinh. Cả ba chính sách đạt 0\% trên hai tập, thỏa ngưỡng 10\%.
\end{itemize}

Abstention F1 được tính bằng $2TP/(2TP+FP+FN)$, với TP là số case thiếu bằng chứng được từ chối đúng. Lỗi sinh ở case thiếu bằng chứng được tính vào FN; lỗi sinh hoặc từ chối trên case có thể trả lời được tính vào FP. Mỗi chính sách có một lỗi sinh ở nhóm thiếu bằng chứng của mỗi tập, tương ứng tỷ lệ sinh thành công 99,29\% và 98,33\%.
\begin{table}[H]\centering\small\setstretch{1.05}
\caption{Chính sách từ chối trên development 140 case và final 60 case.}\label{tab:abstention-results}
\begin{tabular}{@{}llrrrr@{}}\toprule
\textbf{Tập} & \textbf{Chính sách} & \textbf{Abst. F1} & \textbf{False-answer} & \textbf{Token F1} & \textbf{Cit. Valid.}\\\midrule
Development & B0 & 0,9740 & 5,06\% & 0,4490 & 1,0000\\
 & B1 & 0,9936 & 1,27\% & 0,3420 & 1,0000\\
 & B2 & 0,9936 & 1,27\% & 0,3420 & 1,0000\\\midrule
Final & B0 & 0,9697 & 5,88\% & 0,4766 & 1,0000\\
 & B1 & 0,9538 & 8,82\% & 0,2996 & 1,0000\\
 & B2 & 0,9538 & 8,82\% & 0,2996 & 1,0000\\\bottomrule
\end{tabular}
\end{table}

\begin{itemize}
\item \textbf{B1 trên development:} false-answer giảm từ 5,06\% xuống 1,27\%, nhưng Token F1 của 61 case trả lời được giảm khoảng 0,1070 (CI95 $[-0{,}1563;\,-0{,}0631]$), vượt guardrail 0,02.
\item \textbf{B2:} ngưỡng reranker thô khoảng $-7{,}8398$, chọn từ 107 ngưỡng. Cổng không chặn case nào ở development; ở final chỉ chặn một case mà B1 đã từ chối, nên không đổi quyết định nào.
\item \textbf{Quyết định và final:} giữ B0 theo development. Trên final, B0 có 2/34 lỗi ở nhóm thiếu bằng chứng, B1/B2 có 3/34.
\end{itemize}

\textbf{Giới hạn.} B1 đồng thời đổi JSON và chỉ dẫn về hình thức đáp án; mức giảm Token F1 không cô lập tác động của JSON và chưa đủ để kết luận sai ngữ nghĩa. Kết luận chỉ áp dụng cho bộ case đã kiểm soát và các chính sách được thử.

